% Tapis de Quillen, with Quillen's offprint and preprint (cote 162-5, typed part): archive cote 162-5, site page 16 (PDF page 17).
% typescript. Draft compared: none.
% AI-transcribed by Claude Opus 5.5 (claude-opus-5-5, adaptive thinking, effort high), Message Batches API, batch t162_full (msgbatch_01UAh3zyWPd454AeDkhnmQ33), 2026-11-13; not yet checked by eye.
\scanpage{162-5/16}
\mspage{15}

% NOTE: the right edge of the typescript is cut off on the scan; line ends are transcribed as far as visible
\deleted{tés, les morphismes \uline{toutes} les co\unsure{r}respondances \inserted[typed]{co}homologiques (en ve
du thorème de Thom rappelé plus haut). On trouve donc un foncteur
pleinement fidèle de cette catégorie dans celle des vectoriels
gradués \xout{dxxd} positivement de dim. finie sur $\underline{F}_2$ , qui est déjà abéli
ne, et si on veut rendre comme plus haut $\underline{F}_2(-1)$ inversible, on trou
tous les vectoriels gradués de dim. finie sur $\underline{F}_2$. C'est un peu tro
trivial comme structure, et on peut faire mieux,} \xout{exprenantxcomme}
\inserted[typed]{Pour trouver une catégorie universelle,}
\xout{correspondancesxPourxceci} on associe à toute X \xout{\illegible{2}}
\xout{lx} l'ensemble \inserted[typed]{S(X)} des sous-variétés \xout{S(X)} de X, modulo équivalence homo
gique. Donc S(X) $\subset$ H$^{\cdot}$(X), et on prouve qu'on trouve en fait un sous
anneau de H$^{\cdot}$(X) (grâce à un théorème de Thom qui caractérise S(X) \hand{en}
termes d'opérations de Steenrod, cf plus bas). On prouve de même
(par le même théorème) que les S(XxY) sont stables par composition
correspondances cohomologiques( cela ne semble pas trivial, puisque
les S(X) ne sont pas stables par Gysin...). D'où une catégorie add
tive dont les objets sont ceux de C, les flêches sont les \inserted{él. des} S(XxY),
Il semble qu'elle doive jouer le rôle universel que j'ai envisagé
plus haut. En tous cas, sa définition géométrique est fort simple
Nous allons en donner une interpétation algébrique. Pour ceci,
considérons l'anneau des opérations cohomologiques \inserted[typed]{``stables''} \deleted{additives sur l}
\deleted{des endomorphismes du}
\margin{\unsure{?}}\xout{lx} \deleted{foncteur X --- H$^{\cdot}$(X) qui sont additifs, et qui en degré zéro se
réduisent à un multiple de l'identité}\hand{/} C'est un anneau gradué, avec
application diagonale anticommutative, \inserted[typed]{i.e. commutative (car. 2 !)} associative et unitaire. Son
dual gradué est une algèbre \inserted[typed]{graduée anti\xout{o}commutative donc \xout{commutative}} commutative,avec application diagonale
associative et unitaire; il est réduit à $\underline{F}_2$ en degré zéro. Donc il
corespond à un schéma en groupes \unsure{affine} G sur $\underline{F}_2$ , qui est d'aille
le foncteur des \unsure{$\otimes$}-automorphismes \inserted[typed]{additifs} du foncteur H$^{\cdot}$ sur C \hand{qui opèrent
trivialement sur \deleted{\illegible{1}} \unsure{$H^2(-1)$|$\underline{F}_2(-1)$}. \deleted{\illegible{1}}}

\margin{16}
